% =====================================================================
%  Paper I  --  KHI en RRMHD  --  ESQUELETO The Open Journal of Astrophysics
%
%  openjournal.cls se descargo de la propia revista:
%    http://www.thphys.nuim.ie/staff/pcoles/openjournal.cls   (v. 6-sep-2015)
%  Esta basada en emulateapj (estilo ApJ/ApJL) y requiere revtex4-1.
%
%  COMO COMPILAR:  ./build.sh
%
%  ANTES DE ESCRIBIR, LEER:
%    ../manuales/OJAP_instrucciones.md    requisitos de la revista
%    ../02_PLAN_EXTRACCION.md             que sale de donde en la monografia
%    ../04_ESTILO_INGLES.md               ingles + glosario ES->EN
%    ../05_CHECKLIST_ENVIO.md             checklist antes de enviar
%
%  RUTA DE ENVIO: primero arXiv (licencia CC-BY, categoria astro-ph.HE,
%  cross-list astro-ph.IM), y luego "Submit a manuscript" en astro.theoj.org.
%
%  Opciones de clase utiles: [twocolumn] [onecolumn] [apj] [apjl] [tighten]
%                            [numberedappendix] [appendixfloats] [twocolappendix]
% =====================================================================

\documentclass[apj]{openjournal}

\usepackage{graphicx}
\usepackage{amsmath}
\usepackage{amssymb}

\graphicspath{{figures/}{../../monografia/figuras/}}

% --- Macros del proyecto ------------------------------------------------
\newcommand{\gKHI}{\gamma_{\mathrm{KHI}}}
\newcommand{\scrit}{\sigma_{\mathrm{crit}}}
\newcommand{\Rmst}{\mathrm{Rm}^{*}}
\newcommand{\Ozp}{\Omega_{zp}}
\newcommand{\Omaxzp}{\Omega_{zp}^{\max}}
\newcommand{\akh}{a_{\mathrm{kh}}}
\newcommand{\vsh}{v_{\mathrm{sh}}}
\newcommand{\B}{\mathbf{B}}
\newcommand{\kvec}{\mathbf{k}}
\newcommand{\todo}[1]{\textbf{[TODO: #1]}}

% =====================================================================
\shorttitle{Resistive KHI in RRMHD}
\shortauthors{Martinez Anzola et al.}

\begin{document}

\title{\todo{titulo completo} \\
  % Propuesta: The resistive-to-ideal transition of the relativistic
  % Kelvin--Helmholtz instability
}

\author{B. Martinez Anzola}
\author{S. Rodriguez Garcia}
\author{S. Miranda-Aranguren}
\affil{Programa Academico de Fisica, Facultad de Ciencias Matematicas y
  Naturales, Universidad Distrital Francisco Jose de Caldas, Bogota, Colombia}
% TODO: confirmar afiliaciones, orden de autoria y correo de correspondencia.

\begin{abstract}
\todo{Un parrafo. OJAp no publica un limite duro, pero conviene mantenerlo en
250 palabras para que sirva tambien para MNRAS y para arXiv. Contenido: barrido
de 44 simulaciones RRMHD en $\sigma\in[100,10500]$ con el codigo \textit{Cueva};
modelo sigmoide con \emph{offset} y asintota ideal $C+\gamma_0=1.03\pm0.05$
(8 per cent de la teoria lineal RMHD compresible); la transicion es una ZONA
$\sigma\in[1400,7000]$, centro $\scrit\approx2815$; ley
$\Omaxzp\propto\sigma^{1.22}$ por encima del piso Sweet--Parker $\sigma^{1/2}$;
una frase sobre la campana magnetica. Escribirlo AL FINAL.}
\end{abstract}

\keywords{\todo{En OJAp mandan las categorias arXiv: astro-ph.HE (primaria),
astro-ph.IM (cross-list). Anadir aqui palabras clave libres:
magnetohydrodynamics (MHD) --- instabilities --- magnetic reconnection ---
relativistic processes --- methods: numerical}}

% =====================================================================
\section{Introduction}
\label{sec:intro}
\todo{Ver el esqueleto gemelo ../mnras/main.tex y ../02\_PLAN\_EXTRACCION.md
seccion 1. El contenido es identico entre las dos versiones; lo unico que
cambia es la clase, el estilo de citas y (para MNRAS) el ingles britanico.}

\section{Numerical method and setup}
\label{sec:method}
\subsection{Governing equations}
\todo{../02\_PLAN\_EXTRACCION.md seccion 2.1}
\subsection{The \textit{Cueva} code}
\todo{../02\_PLAN\_EXTRACCION.md seccion 2.2}
\subsection{Initial conditions}
\todo{../02\_PLAN\_EXTRACCION.md seccion 2.3}
\subsection{The two campaigns}
\todo{../02\_PLAN\_EXTRACCION.md seccion 2.4}

\section{Results: the conductivity sweep}
\label{sec:results}
\todo{../02\_PLAN\_EXTRACCION.md seccion 3 -- los tres resultados centrales.}

\section{Comparison with linear theory}
\label{sec:theory}
\todo{../02\_PLAN\_EXTRACCION.md seccion 4}

\section{Magnetic field strength and orientation}
\label{sec:magnetic}
\todo{../02\_PLAN\_EXTRACCION.md seccion 5}

\section{Discussion and conclusions}
\label{sec:conclusions}
\todo{../02\_PLAN\_EXTRACCION.md seccion 6}

% =====================================================================
\acknowledgments

\todo{Recursos de computo, financiacion, personas, software citado.
IMPORTANTE -- OJAp exige que la DECLARACION DE USO DE IA vaya aqui, en los
agradecimientos, y que si se uso IA para calculo, analisis o visualizacion se
describan los pasos de validacion. Ver ../comun/DECLARACION\_IA.md.}

% OJAp no exige una seccion "Data availability", pero recomienda seguir
% Chen et al. (2022, ApJS 260, 5). La incluimos igual: ver
% ../comun/DATA_AVAILABILITY.md
\section*{Data availability}
\todo{Plantilla A de ../comun/DATA\_AVAILABILITY.md}

% =====================================================================
% aasjournalv7.bst viene con aastex en TeX Live y da el estilo ApJ autor-anio.
\bibliographystyle{aasjournalv7}
\bibliography{references}
% Ver ../comun/NOTAS_BIBTEX.md (conversion biblatex -> natbib).

\appendix
\section{Supplementary statistical detail}
\todo{Opcional; ver ../02\_PLAN\_EXTRACCION.md seccion 8.}

\end{document}
